% Chapter 4 -- translation by Doug Burkett & Tim Hutcheson (ca. 2001),
% lightly copy-edited for this edition.

\chapter{Expert simulation in direct current}

\section{Introduction to the (Expert) Mode}

If you are not interested in learning to simulate like an expert,
you can skip this chapter. In this chapter, I describe advanced
concepts of Symbulator operation, as well as examples that show how
you can use the (Expert) Mode to save time during symbolic
simulations. To get the most benefit, you must read this chapter
attentively and thoroughly.

In the previous chapter we did several symbolic simulations, and we
learned to use the \code{solve} command and the \code{solves} tool.
In this chapter, we will learn to simulate with the (Expert) Mode.
This mode appeared in version 3 of the Symbulator, in order to give
the user more control over the simulation process.

In the simulations we have seen so far, the Symbulator behaves like
a black box, in which you supply the circuit description and
Symbulator gives you the results. And that is all: while you are
involved in complimenting the girl sitting next to you, or
contemplating the immortality of the crab, Symbulator is working to
generate the equations, solve them and save the answers. Your only
responsibility is to correctly enter the circuit description, and
Symbulator does the rest.

The (Expert) Mode gives you the power to play a more important role
in the simulation, because it gives you the capability to modify the
equations which Symbulator generates, before the calculator solves
them. In return, you assume part of the responsibility for the
simulation, which is why the (Expert) Mode must be used with care.
To use Symbulator as an expert requires deeper knowledge than is
required of the casual user. Nevertheless, for those willing to
learn a little more, the (Expert) Mode is an attractive option,
because it offers faster simulations.

\section{Theory}

\subsection{The procedure in general}

An expert simulation, as the word implies, requires that you be an
expert, not only in the use of Symbulator, but also an expert in its
internal operations. Let us learn, then, the internal simulation
processes. In general, a simulation done by Symbulator consists of
these steps:

\begin{enumerate}
\item Equations are generated to describe the circuit.
\item These equations are solved.
\item The answers are stored.
\end{enumerate}

Within this general framework, we can say that the (Expert) Mode
allows us to manipulate the equations generated in step 1, before
they are solved in step 2, so that when the answers are stored in
step 3, they are in the form we want.

\subsection{First-level equations}

The necessary equations used by Symbulator to solve a circuit are
called equations of the first level. Symbulator generates two types
of first-level equations: node equations, and special equations.

There is a node equation for each circuit node. These equations are
generated using Ohm's law and Kirchhoff's current law. The
Symbulator will generate for each node an equation which fulfills
Kirchhoff's current law: the sum of all the currents of the elements
in contact with that node, set equal to zero. To describe the
currents of resistive elements, Symbulator uses Ohm's law.

In addition, a special equation is generated for each special
element in the circuit. These equations are generated according to
the theoretical formulas which describe the behavior of the
elements. It may seem strange, but voltage sources and short
circuits are special elements. This results because Symbulator uses
nodal analysis, so the only elements (seen so far) which are not
special are resistors and current sources. The others are considered
special, and each is associated with a special equation, which will
increase the number of equations generated by the Symbulator.

For example, to solve a circuit with five nodes, two voltage sources
and one short circuit, Symbulator will generate a total of eight
first-level equations. It does not matter if this circuit has one
thousand resistors and seven hundred current sources, it will still
have eight first-level equations. However, the complexity of the
eight equations increases as more resistors and current sources are
connected to the nodes.

How many first-level equations will be generated to simulate a
circuit with three nodes and one voltage source? Four: one for each
node, and one more for the voltage source.

In our algebra courses, we learned that we need $n$ equations to
completely solve a system of $n$ unknowns. This means that for the
first circuit, which has eight equations, we can solve for eight
unknowns. In the second circuit, we can solve for four unknowns,
since we have four equations. The unknowns which can be solved with
the first-level equations are called first-level variables.

\subsection{First-level variables}

As expected, Symbulator generates the first-level equations based on
the first-level variables. There are as many first-level variables
as there are first-level equations. Which means, indirectly, that
there will be as many first-level variables as there are nodes,
voltage sources and short circuits in the circuit.

With the elements we have seen so far, Symbulator generates three
types of first-level variables: node voltages, voltage source
currents, and short circuit currents.

Each node voltage in the circuit is a first-level variable. When
solving a circuit which has three nodes 1, 2 and 3, Symbulator
creates three first-level variables called \code{v1}, \code{v2} and
\code{v3}.

The current through each voltage source is a first-level variable.
When solving a circuit with two voltage sources \code{e1} and
\code{ex}, Symbulator creates two first-level variables called
\code{ie1} and \code{iex}.

In addition, each short-circuit current is a first-level variable.
So, when solving a circuit with a short circuit called \code{s1},
Symbulator creates a first-level variable called \code{is1}.

After solving the first-level equations, the results are stored in
the first-level variables. Can you guess what these answers are
called? Obviously: first-level answers. In the example circuit with
five nodes, two voltage sources and one short circuit, eight
first-level equations are created based on the eight first-level
variables. The simulation results are the five node voltages, the
two voltage-source currents, and the single short-circuit current.

How many first-level variables will Symbulator generate for the
circuit with three nodes and one voltage source? Obviously, there
will be four: one for the voltage source current, and three for the
node voltages. After the four first-level equations are solved,
there are four first-level results.

\subsection{Second-level variables}

I have described currents in voltage sources and short circuits as
first-level variables. However, currents in other elements such as
current sources and resistors are called second-level variables.

Voltage drops across resistors and sources are also second-level
variables. So, in a circuit with four nodes, three resistors, two
current sources, two voltage sources, and one short circuit, we will
have these first-level variables:

\begin{itemize}
\item 4 node voltages
\item 2 currents in the voltage sources
\item 1 current in the short circuit
\end{itemize}

and these second-level variables:

\begin{itemize}
\item 3 currents in the resistors
\item 2 currents in the current sources
\item 7 voltage drops across the 3 resistors and the 4 sources
\end{itemize}

This results in a total of $4 + 2 + 1 = 7$ first-level variables,
with their respective first-level equations, and $3 + 2 + 7 = 12$
second-level variables. Note that the short circuit does not
generate a voltage drop (see section 2.6.3 on short circuit
results).

The second-level variables are not included in the first-level
equations. So, how are these variables related to the simulation
results?

\subsection{Second-level expressions}

The answer is simple. For each second-level variable, the Symbulator
generates a second-level expression, which is a function only of the
first-level variables. So, given the first-level results, all the
second-level variables are immediately defined by these expressions.

There are two types of second-level expressions: those which are
replaced by their results after the solution, and those which remain
expressed in terms of the first-level variables. That is, after the
first-level equations are solved and the results are stored in the
first-level variables, Symbulator evaluates some of the second-level
expressions, and so generates second-level results which are stored
in the appropriate variables. However, other second-level
expressions are not evaluated and remain expressions instead of
results. Second-level expressions which are not evaluated are those
that correspond to the currents in current sources, and the voltage
drops. These remain algebraic expressions as functions of the node
voltages.

\subsection{Third-level variables}

In direct current (DC) and alternating current (AC) simulations,
there is another group of variables called third-level variables.
They are given this name because they are generated only after the
solution of the first-level equations and the evaluation of the
second-level expressions. In other words, they are only created when
the previous levels of variables have both been solved.

The third-level variables are the powers consumed by the elements.
They are of course generated according to the definition of
electrical power, which is the product of the voltage drop across an
element and the current through the element, which must be
previously defined.

\subsection{The normal procedure in detail}

We have learned about first-level equations, variables and results,
second-level expressions and variables, and third-level variables. I
will now describe the simulation process in detail, as it relates to
these variables. First, the procedure of a normal simulation, which
does not use the (Expert) Mode:

\begin{enumerate}
\item First-level equations are generated which describe the circuit
  entered by the user, based on the first-level variables. The
  second-level expressions are also generated, as functions of the
  first-level variables.
\item The first-level equations are solved for the first-level
  variables.
\item The first-level results are stored in the first-level
  variables. Some second-level expressions are evaluated, and their
  second-level results are stored.
\item In direct current and alternating current analyses, the
  third-level variables are generated, which are the powers consumed
  by the elements.
\end{enumerate}

\subsection{Procedure options for the (Expert) Mode}

Several options are available when simulating with the (Expert)
Mode. The first option is to choose the correct analysis type. These
options are direct current (DC) analysis, alternating current (AC)
analysis, frequency domain (FD) analysis and transient analysis
(TR).

Choose the desired analysis type and press \key{ENTER}. Next, the
calculator shows a form from which you choose two options that
determine how floating-point numbers are handled:

The first option, Digits?, sets the number of digits to be used in
the calculations. Normally this value should be set to 9, unless a
different precision is required for some reason.

The second option, Format?, is used to set the display format. You
can set the display format to Normal, Engineering or Scientific.

Press \key{ENTER} after you have chosen the desired options. At this
point, Symbulator executes the first step of the simulation,
detailed above. Before the second step, the program pauses and shows
another form, which offers you the opportunity to manipulate the
first-level equations and unknowns. You can change them or not, as
needed.

This screen is called ``(Expert) Mode -- 1st Level''. There are four
fields and a menu of options:

The first field shows the first-level unknowns. If there is more
than one unknown, they are shown separated with commas between curly
brackets: \{\}.

The second field shows the first-level equations. If there is more
than one equation, they are shown connected by the \code{and}
operator.

The third and fourth fields are used to specify known first-level
variables and values. The third field is used to specify the name of
a first-level variable whose value is known. These variables are
called conditional variables. If there is more than one known
variable, they must be separated with commas, for example,
\code{name1,name2,name3}. The fourth field is used to specify the
values of the variables listed in the third field. These declared
values are called conditional equations. If there is a single
variable, its value is specified as an equality, for example,
\code{v1=3.5}. If there is more than one, the equalities are
connected with the \code{and} operator, for example, \code{v1=3.5
and v2=5 and ie1=-2}. Note that only first-level variables can be
specified in these fields.

The menu shown under the four fields has three options which let you
control the simulation process. The options are 1) run the
simulation, 2) run the simulation and keep the equations, and 3)
only store the equations.

Let us see next what each menu option means.

\subsection{Only store the equations}

If you choose the third option, ``Just Keep'', then Symbulator will
simply create four variables:

\begin{itemize}
\item \code{Unknown}, which holds the list of the first-level
  unknowns
\item \code{Equation}, which holds the first-level equations
\item \code{Wheneq}, which holds the conditional equations
\item \code{Whenvar}, which holds the conditional variables
\end{itemize}

Each variable is stored as a text string. Any modifications the user
may have introduced while they were displayed in the dialog box are
reflected in these variables. The purpose of the ``Just Keep''
option is to keep the equations and other variables so that the
\code{solves} tool can be used later. If the circuit is very large,
the generated equations may be too large to be shown on the screen.
In this case, the equations are not displayed, but instead a message
is shown indicating that the equations will be stored automatically.
However, in most cases the equations will fit on the display screen.

After the four variables are created, Symbulator shows a new screen
called ``(Expert) Mode -- 2nd Level''. You use this screen to choose
whether or not you want to keep the second-level expressions which,
as we have said, have already been generated.

You choose to either keep the expressions or delete them. If you
choose to delete them, Symbulator deletes them and finishes its
work. If you choose to keep them, they are stored, and a third
screen is displayed, called ``(Expert) Mode -- 3rd Level''. In this
screen you choose whether or not to generate the third-level
expressions.

If you choose not to generate the expressions, Symbulator exits. If
you choose to generate the expressions, Symbulator does that, then
exits.

\subsection{Simulate and store}

If you choose the second option, ``Symbulate+Keep'', that is, to
simulate and keep the equations, then Symbulator creates the four
variables described in the previous section, then the simulation
proceeds as usual. That is, steps 2, 3 and 4 are executed as they
would be normally, taking into account the changes you have made.

\subsection{Simulate only}

If you choose the first option, ``Symbulate'', then Symbulator runs
the simulation without creating the four variables. Steps 2, 3 and 4
are executed, taking into account the changes you have made,
obviously. Now that we have learned the theory, let us see all these
new concepts applied in practice. We will solve the problems of the
previous chapter, but this time, like experts.

\section{Practice}

\problemx{013}{with the (Expert) Mode}

\emph{Problem: Determine the numerical values of $v_x$ and $i_x$ in
the following circuit.}

Solution:

\subsection{Initial steps}

We name the nodes and elements, and delete the symbolic names:

\begin{entry}
DelVar ra,vx
\end{entry}

We enter the circuit description and specify expert simulation:

\begin{entry}
sq\expert("e1,1,0,18;ra,1,0,ra;r1,1,0,6;r2,2,1,5;ex,2,0,vx")
\end{entry}

We start the simulation with \key{ENTER}. We choose DC analysis and
press \key{ENTER}. We leave the floating-point formats unchanged and
press \key{ENTER}. The screen shows status messages indicating that
the (Expert) Mode is activated and that the first step is executed.
The (Expert) Mode 1st Level screen appears. In this screen we can
see the equations generated by Symbulator, and the unknowns for
which it intends to solve them.

Remember that resistor \code{r2} was defined so that the direction
of its current agrees with that of the 12~A current given in the
problem as a known value. However, if we read the list of
first-level unknowns in the first field, we see that this resistor
current is not shown, for it is not a first-level unknown. On the
other hand, the currents of the two voltage sources, \code{ie1} and
\code{iex}, are shown, for they are first-level unknowns.

Since the current of the \code{ex} source is defined from node 2 to
node 0, we can say that the current \code{iex} is equal to $-12$~A.
The negative sign is used since the source current is defined
opposite the direction shown in the problem.

Now the time has come to think carefully about the problem, and what
results are needed. We see that Symbulator has generated a system of
four equations and four unknowns, which also has one symbolic value.
We consider this symbolic value as an unknown, because it must be
solved to obtain a numerical answer. So, we really have a system of
four equations and five unknowns. We can look at this situation from
two perspectives: we either have an extra unknown, or we need
another equation.

We will solve this problem first from the first perspective, and
then from the second.

\subsection{Eliminating an unknown}

If we have an extra unknown, we must eliminate it so we can solve a
system of four equations and four unknowns. We take three steps to
eliminate the unknown. First, we delete \code{iex} from the list of
unknowns in the first field. We do not need it, because its value
will be declared as known. Second, we add the new first-level
variable \code{vx} to the list of unknowns. Finally, we add the name
\code{iex} to the third field, and the equality \code{iex=-12} to
the fourth field, where the `$-$' sign is entered with \key{($-$)}.
After making these modifications, the fields look like this:

\begin{quote}
First-level unknowns: \code{\{v1,ie1,v2,vx\}}\\
The first-level equations were not modified.\\
Known variable names: \code{iex}\\
Known variable values: \code{iex=-12}
\end{quote}

We press \key{ENTER} to continue the simulation, and the screen
shows the status messages for the remaining three simulation steps,
then `Done' is shown.

Unlike the previous techniques which require use of the \code{solve}
command and the \code{solves} tool, this (Expert) Mode simulation
does not require additional manipulation to obtain numerical
answers. By modifying the list of unknowns and specifying the known
value, we reduced a system of four equations and five unknowns to a
system of four equations and four unknowns. So, the simulation
results, which are now in the current folder, are numerical results.

By recalling \code{ir1}, we find that the current $i_x$ is 3~A. To
find the value of $v_x$, we have three options: 1) recall the node
voltage \code{v2}, 2) recall the voltage drop \code{vex} across the
source, or 3) simply recall our first-level unknown \code{vx}.
Remember that we included \code{vx} in the list of first-level
unknowns, so Symbulator solved for it as it would for any other
first-level unknown. By any of these three roads, we find that
$v_x = 78$~V.

Now let us use the other perspective to solve the problem.

\subsection{Adding an equation}

If we need another equation, we must add it, to have a system of
five equations and five unknowns. Since we want to do this problem
again from this new perspective, we execute all the steps of section
4.3.1 until the 1st Level screen is shown:

\begin{entry}
DelVar ra,vx
sq\expert("e1,1,0,18;ra,1,0,ra;r1,1,0,6;r2,2,1,5;ex,2,0,vx")
\end{entry}

We start the simulation with \key{ENTER}, choose DC analysis, leave
the floating-point formats as they are, and press \key{ENTER}. The
1st Level screen is shown.

To add an equation we must do two things: 1) add the variable
\code{vx} to the list of first-level unknowns, and 2) add an
equation to the list of first-level equations, like this: \code{and
iex=-12}. After these changes, the field contents are:

\begin{quote}
First-level unknowns: \code{\{v1,ie1,v2,iex,vx\}}\\
First-level equations: \code{(5*ie1*ra-ra*(vx-33)+90)/ra=0 and
5*iex+vx=18 and v1=18 and v2=vx and iex=-12}\\
Known variable names: empty.\\
Known variable values: empty.
\end{quote}

We have changed a system of four equations and five unknowns into a
system of five equations and five unknowns. We press \key{ENTER} to
continue the simulation, the status messages are shown, then `Done'
is shown, indicating that the simulation is complete. The numerical
results are in the current folder.

We recall the value of \code{ir1} to find that $i_x$ is 3~A. We
recall the value of \code{v2}, or \code{vx}, or \code{vex}, to find
that $v_x$ is 78~V. This problem illustrates that the (Expert) Mode
can save an expert user time, but it requires good knowledge of
circuits, algebra and the use of Symbulator.

Now let us see another problem of the previous chapter.

\problemx{014}{with the (Expert) Mode}

\emph{Problem: Determine the numerical values of $v_x$ and $i_x$ in
the following circuit.}

Solution:

We name the nodes and elements, and clear the symbolic names:

\begin{entry}
DelVar rx
\end{entry}

We enter the circuit description and the expert simulation command:

\begin{entry}
sq\expert("j1,0,1,6;r1,1,2,1;r2,1,3,5;j2,3,2,10;r3,3,0,2;
          rx,2,4,rx;r4,4,0,3")
\end{entry}

We start the simulation with \key{ENTER} and choose DC analysis. We
leave the floating-point formats unchanged and press \key{ENTER}.
The display shows the status messages indicating that the (Expert)
Mode is active and that the first step is complete. Then, the 1st
Level screen is shown.

We have a system of four equations and five unknowns, where the
fifth unknown is the symbolic value of \code{rx}, which must be
solved to get numeric values. We will solve this problem by adding
an equation to the system, which is the simplest method. First, we
add \code{rx} to the list of first-level unknowns, so that
Symbulator now considers it an unknown. Second, we add an equation
to the first-level equations list: \code{and ir3=4}. So now the
field contents are:

\begin{quote}
First-level unknowns: \code{\{v1,v2,v3,v4,rx\}}\\
First-level equations: \code{(rx*(v1-v2+10)-v2+v4)/rx=0 and
(rx*v4-3*(v2-v4))/rx=0 and 6*v1-5*v2-v3-30=0 and 2*v1-7*v3-100=0
and ir3=4}\\
Known variable names: empty.\\
Known variable values: empty.
\end{quote}

Now I must explain what is allowed in the (Expert) Mode, and what is
not. Note that \code{ir3} is included in the equations, but is
\code{ir3} a first-level variable? Of course not. It is a
second-level variable. But at this point, step 1 of the simulation
has already been executed, in which the second-level expressions
have been generated, which express the second-level variables as
functions of the first-level variables. Therefore, we are allowed to
include second-level variables in the first-level equations, because
the second-level variables are expressed internally in terms of the
first-level unknowns. On the other hand, you cannot specify
second-level variables with known values in the third and fourth
fields: these fields can only be used for the values of first-level
variables.

Let us return to the problem. After we change the system of four
equations and five unknowns to a system of five equations and five
unknowns, we press \key{ENTER} and the screen messages show us that
the three remaining simulation steps are executed.

Finally the message `Done' is shown, and the numerical answers are
in the current folder. We recall the value of \code{ir1} and find
that $i_x = -8$~A. We recall the value of \code{vrx} and find that
$v_x = 80$~V. It is possible to solve this problem with the other
technique, in which we eliminate an unknown, but that method is a
little more complicated here. We would have to declare the value of
some first-level variable as known, in the third and fourth fields.
In order to express the known current of 4~A through a first-level
variable, we would use Ohm's law to relate the voltage of node 3 and
the 2 ohm resistance: $v_3 = (4~\mathrm{A})(2~\Omega)$, that is,
$v_3 = 8$. The procedure, then, would be:

\begin{entry}
DelVar rx
sq\expert("j1,0,1,6;r1,1,2,1;r2,1,3,5;j2,3,2,10;r3,3,0,2;
          rx,2,4,rx;r4,4,0,3")
\end{entry}

We start the simulation by pressing \key{ENTER}, choose DC analysis,
and leave the floating-point formats unchanged, then press
\key{ENTER} to continue. The display shows that the (Expert) Mode is
activated and the first step is executed, then shows the 1st Level
screen.

To eliminate one system unknown, we must do three things. First, we
remove \code{v3} from the list of first-level unknowns in the first
field, because we know its value. Second, we add \code{rx} to the
list of first-level unknowns. Third, we add the name \code{v3} to
the third field, and specify its value in the fourth field as
\code{v3=8}. After these changes, the fields are:

\begin{quote}
First-level unknowns: \code{\{v1,v2,rx,v4\}}\\
The first-level equations were not modified.\\
Known variable names: \code{v3}\\
Known variable values: \code{v3=8}
\end{quote}

We press \key{ENTER} to continue the simulation, the screen shows
the status messages for the three remaining simulation steps, and
finally `Done' is shown.

We recall the value of \code{ir1}, and find that $i_x = -8$~A. We
recall the value of \code{vrx}, and find that $v_x = 80$~V. This
example shows that when it is difficult to declare a known value
through a first-level variable, the method of adding an equation is
easier than eliminating an unknown. Nevertheless, eliminating an
unknown has the advantage of reducing the system of equations,
making the system easier for the calculator to solve. This is
important in circuits much larger than we see in these pages.

\problemx{015}{with the (Expert) Mode}

\emph{Problem: Determine the numerical values of $v_x$ and $i_x$ in
the following circuit.}

Solution:

We name the nodes and elements, and delete the name \code{ix}:

\begin{entry}
DelVar ix
\end{entry}

We enter the circuit description and specify expert simulation:

\begin{entry}
sq\expert("e1,1,0,60;r1,1,2,8;r2,2,0,10;r3,2,3,4;r4,3,0,2;
          jx,0,3,ix")
\end{entry}

We start the simulation by pressing \key{ENTER}, choose DC analysis,
and leave the floating-point formats unchanged. After pressing
\key{ENTER}, the display shows that the (Expert) Mode is active, and
the first simulation step is executed. The 1st Level screen is
shown. Like the previous problems, we can solve this one by adding
an equation or by eliminating an unknown.

To add an equation, we change the field contents so that we have
five equations and five unknowns:

\begin{quote}
First-level unknowns: \code{\{v1,ie1,v2,v3,ix\}}\\
First-level equations: \code{8*ie1-v2+60=0 and 4*ix+v2-3*v3=0 and
v1=60 and 19*v2-10*(v3+30)=0 and ir1=5}\\
Known variable names: empty.\\
Known variable values: empty.
\end{quote}

On the other hand, to solve the problem by eliminating an unknown,
we change the fields so that we have a system of four equations and
four unknowns, by declaring the value of a first-level variable,
like this:

\begin{quote}
First-level unknowns: \code{\{v1,v2,v3,ix\}}\\
The first-level equations were not modified.\\
Known variable names: \code{ie1}\\
Known variable values: \code{ie1=-5}
\end{quote}

You can use either method, as you prefer. Next, we press \key{ENTER}
and the display shows that the three remaining steps are executed,
then the display shows `Done'.

We recall the value of \code{ix} and find that $i_x = 1$~A. We
recall the value of \code{-vjx}, or of \code{v3}, and find that
$v_x = 8$~V. Let us see one more problem.

\problemx{016}{with the (Expert) Mode}

\emph{Problem: Determine the value of $k$ such that voltage $v_y$ is
zero.}

Solution:

We name the nodes and elements and delete the variable name
\code{k}:

\begin{entry}
DelVar k
\end{entry}

We enter the circuit description and specify expert simulation:

\begin{entry}
sq\expert("e1,1,0,6;r1,1,x,1;r2,x,0,4;r3,x,y,2;j1,0,y,2;
          r4,y,2,3;e2,2,0,k*vx")
\end{entry}

We start the simulation by pressing \key{ENTER}, choose DC analysis,
and leave the floating-point formats unchanged. The display shows
that the (Expert) Mode is active and that the first simulation step
is complete. Next the display shows the 1st Level screen. This
problem can be solved by adding an equation or by eliminating an
unknown.

To add an equation, the field contents are modified so that we have
seven equations and seven unknowns:

\begin{quote}
First-level unknowns: \code{\{v1,ie1,vx,vy,v2,ie2,k\}}\\
First-level equations: \code{2*k*vx+3*vx-5*vy+12=0 and
k*vx+3*ie2-vy=0 and ie1-vx+6=0 and v1=6 and v2=k*vx and
7*vx-2*(vy+12)=0 and vy=0}\\
Known variable names: empty.\\
Known variable values: empty.
\end{quote}

Alternatively, to solve the problem by eliminating an unknown, we
modify the field contents so we have six equations and six unknowns,
by declaring the value of a first-level variable, like this:

\begin{quote}
First-level unknowns: \code{\{v1,ie1,vx,v2,ie2,k\}}\\
The first-level equations were not modified.\\
Known variable names: \code{vy}\\
Known variable values: \code{vy=0}
\end{quote}

You can use either method, as you prefer. Next we press \key{ENTER}
and the display shows the execution of the three remaining
simulation steps, and finally `Done' is shown. We recall the value
of \code{k} and find that it is $-13/4$.

\section{The solves tool after expert simulation}

The \code{solves} tool can be used at any time. In the previous
chapter we used it after a symbolic simulation. In that case, the
variables \code{Equation} and \code{Unknown} did not exist in the
current folder. When we used \code{solves} for the first time, the
fields were empty, and we manually entered the equations we wanted
to solve.

When we next used \code{solves}, the fields showed the previous
manually-entered equations, unless the folder contents had been
deleted.

To this point we have done our expert simulations using the first
option of the (Expert) Mode, which is to simulate only, without
keeping the equations. For that reason, the variables
\code{Equation} and \code{Unknown} are not created in the current
folder.

On the other hand, if we had used either of the other two options,
that is, simulate and keep the equations, or only keep the
equations, then we would have the variables \code{Equation} and
\code{Unknown} in the current folder.

The \code{solves} tool works differently when executed just after an
expert simulation which has kept the generated equations. In this
case, the fields are not empty, but instead show the equations and
unknowns which were stored by the (Expert) Mode in the variables
\code{Equation} and \code{Unknown}.

For example, if we simulate a circuit with the (Expert) Mode and
choose the ``Just Keep'' option, then Symbulator will store the
equations and unknowns, and we can later solve them with the
\code{solves} tool. Note that if you intend to insert a second-level
variable in the first-level equations, you must ask the (Expert)
Mode to keep the second-level expressions which Symbulator
generated.

The \code{solves} tool shows two menus in its lower part. You use
the first menu to choose whether to solve the equations as real or
complex. You must specify real solutions unless the equations of an
AC analysis are being solved. You use the second menu to choose
whether or not to use conditions. What are the conditions? They are
nothing more than the known values of first-level variables. If we
choose to use conditions, a new screen is shown so that you can
enter these values and their names.

In the field called ``When eq's'', the known values of the
first-level variables must be entered. In the field called ``When
var's'', the names of these first-level variables are entered.

\section{Verification of numerical results of expert simulations}

We can use the procedure of section 3.7 to verify the numerical
results of expert simulations.

Before closing this chapter, I want to congratulate you on deciding
to become an expert Symbulator user. Now that you know the
theoretical concepts, all that remains is to use them in practice
for some time. The rewards for your efforts, in satisfaction and
time saved, will not be long in coming.
